\centerline{\bf \Huge Банки. База(А+Д) II}
\centerline{\bf \Huge }

\begin{flushright}
        \scriptsize{\textbf{qq}}
\end{flushright}


\problem{\textbf{1}} В июле 2025 года планируется взять кредит в банке на сумму 600 тысяч рублей на 10 лет. Условия его возврата таковы:

— каждый январь долг будет возрастать на $r$\% по сравнению с концом предыдущего года;

— с февраля по июнь каждого года необходимо выплатить одним платежом часть долга;

— в июле 2026, 2027, 2028, 2029 и 2030 годов долг должен быть на какую-то одну и ту же
величину меньше долга на июль предыдущего года;

— в июле 2030 года долг должен составлять 400 тыс. рублей;

— в июле 2031, 2032, 2033, 2034 и 2035 годов долг должен быть на другую одну и ту же
величину меньше долга на июль предыдущего года;

— к июлю 2035 года кредит должен быть выплачен полностью.

Известно, что сумма всех платежей после полного погашения кредита будет равна 1740 тыс.
рублей. Найдите $r$.

\problem{\textbf{2}} В июле 2025 года планируется взять кредит на 600 тыс. рублей. Условия его возврата таковы:

— в январе 2026, 2027 и 2028 годов долг возрастает на $r$\% по сравнению с концом предыдущего года;

— в январе 2029, 2030 и 2031 годов долг возрастает на 15\% по сравнению с концом предыдущего года;

— в июле каждого года долг должен быть на одну и ту же величину меньше долга на июль
предыдущего года;

— к июлю 2031 года долг должен быть полностью погашен.

Чему равно $r$, если общая сумма выплат составит 930 тыс. рублей?

\problem{\textbf{3}} 15 января планируется взять кредит в банке на некоторую сумму на 31 месяц. Условия его возврата таковы:

— 1-го числа каждого месяца долг увеличивается на 1\% по сравнению с концом предыдущего
месяца;

— со 2-го по 14-е число каждого месяца необходимо выплатить одним платежом часть долга;

— на 15-ое число каждого месяца с 1-го по 30-й долг должен должен быть на 20 тыс. рублей
меньше долга на 15-е число предыдущего месяца;

— к 15-му числу 31-го месяца долг должен быть погашен полностью.

Сколько тысяч рублей составляет долг на 15 число 30-ого месяца, если банку всего было
выплачено 1348 тыс. рублей?